\documentclass{article}
\usepackage[utf8]{inputenc}
\usepackage[T1]{fontenc}
\usepackage[french]{babel}
\usepackage{amsmath, amsfonts, amssymb}
\usepackage{geometry}
\geometry{a4paper, margin=1in}

\begin{document}

\section*{Énoncé}

\paragraph{Définition}
Soit $f$ une fonction continue sur un intervalle $I$, $a$ et $b$ deux réels de $I$ et $F$ une primitive de $f$ sur $I$.
On appelle intégrale de $f$ entre $a$ et $b$ le réel, noté $\int_a^b f(x)dx$, défini par $\int_a^b f(x)dx = F(b)-F(a)$.
On note $\int_a^b f(x)dx = [F(x)]_a^b = F(b)-F(a)$.

\paragraph{Propriétés}
Soit $f$ une fonction continue sur un intervalle $I$, $a$, $b$ et $c$ des réels de $I$.
\[
\int_a^b \alpha dx = \alpha(b-a), \quad \alpha \in \mathbb{R}; \quad \int_a^a f(x)dx = 0; \quad \int_a^b f(x)dx = -\int_b^a f(x)dx; \quad \int_a^c f(x)dx + \int_c^b f(x)dx = \int_a^b f(x)dx.
\]

\paragraph{Théorème}
Soit $f$ et $g$ deux fonctions continues sur un intervalle $I$, $a$ et $b$ deux réels de $I$.
Pour tous réels $\alpha$ et $\beta$,
\[
\int_a^b (\alpha f(x)+\beta g(x))dx = \alpha \int_a^b f(x)dx + \beta \int_a^b g(x)dx.
\]

\paragraph{Théorème}
Soit $f$ une fonction continue sur un intervalle $I$, $a$ et $b$ deux réels de $I$.
\begin{itemize}
    \item Si $f$ est positive sur $I$ et $a \le b$, alors $\int_a^b f(x)dx \ge 0$.
    \item Si $f$ est négative sur $I$ et $a \le b$, alors $\int_a^b f(x)dx \le 0$.
\end{itemize}

\paragraph{Corollaire}
Soit $f$ une fonction continue sur un intervalle $I$, $a$ et $b$ deux réels de $I$.
\begin{itemize}
    \item Si $\forall x \in I, g(x) \le f(x) \le h(x)$ et $a \le b$, alors $\int_a^b g(x)dx \le \int_a^b f(x)dx \le \int_a^b h(x)dx$.
    \item Si $a \le b$, alors $\left|\int_a^b f(x)dx\right| \le \int_a^b |f(x)|dx$.
\end{itemize}

\paragraph{Théorème}
Soit $f$ une fonction continue sur un intervalle $I$, $a$ et $b$ deux réels de $I$.
Si $f$ est positive et ne s'annule qu'en un nombre fini de réels de $I$ et $a < b$ alors $\int_a^b f(x)dx > 0$.

\paragraph{Théorème}
Soit $u$ et $v$ deux fonctions dérivables sur $[a,b]$ et telles que leurs dérivées $u'$ et $v'$ sont continues sur $[a,b]$. Alors
\[
\int_a^b u(t)v'(t)dt = [u(t)v(t)]_a^b - \int_a^b u'(t)v(t)dt.
\]

\paragraph{Théorème}
Le plan est muni d'un repère orthogonal.
Soit $f$ et $g$ deux fonctions continues sur $[a,b]$.
L'aire (en u.a) de la partie du plan limitée par la courbe de $f$, la courbe de $g$ et les droites d'équations $x=a$ et $x=b$ est le réel $\int_a^b |f(x)-g(x)|dx$.

\end{document}
